\chapter{索引}
\label{app:index}

\section*{索引使用}
\label{app:index-guide}

本索引只链接已经写入正文的锚点、材料页和图件。事件账本仍保存全部 ID；下列七组“核验锚点簇”把其中 1,304 条放回王朝年序，并以行动者、空间、前后链和材料限度作首次描写。余下 73 条分散在各年段，列于账本的延后核验清单，不能借索引伪装成已写正文。

\section*{事件索引}
\label{app:event-index}

\begin{description}
  \item[明初] \eventref{EM-1368-FOUNDING}{1368：朱元璋在应天即皇帝位并建国号明}
  \item[土木堡] \eventref{EM-1449-TUMU}{1449：英宗亲征瓦剌、明军在土木堡溃败}
  \item[荆襄] \eventref{MM-1465-JINGXIANG-UPRISING}{1465：荆襄流民起事}
  \item[大藤峡] \eventref{MM-1465-DATENG-GORGE-CAMPAIGN}{1465：大藤峡征讨}
  \item[建州] \eventref{MM-1467-JIANZHOU-CAMPAIGN}{1467：建州之役}
  \item[宁王之役] \eventref{MM-1519-NING-REBELLION}{1519：南昌的反叛}
  \item[大礼议] \eventref{MM-1522-GREAT-RITES-BEGIN}{1522：大礼议的开端}
  \item[东南海防] \eventref{MM-1547-ZHU-WAN-APPOINTMENT}{1547：朱纨受命整饬浙闽海防}
  \item[月港] \eventref{MM-1567-YUEGANG-OPENING}{1567：月港有限开海}
  \item[隆庆边市] \eventref{MM-1571-LONGQING-PEACE}{1571：封俺答与恢复互市}
  \item[援朝] \eventref{ML-1592-902}{1592：明廷开始援朝调兵}
  \item[播州] \eventref{ML-1606-901}{1606：明军攻破海龙屯、播州改流}
  \item[辽东] \eventref{ML-1618-901}{1618：努尔哈赤向明辽东发动攻势}
  \item[三饷] \eventref{ML-1637-901}{1637：增设练饷}
\end{description}

\section*{王朝年序与行动者索引}
\label{app:polity-index}

\begin{description}
  \item[明中央、官僚与军队] \hyperref[register:1368-1398_hongwu]{洪武}、\hyperref[register:1399-1424_jianwen_yongle]{建文--永乐}、\hyperref[register:1425-1464_hongxi_tianshun]{洪熙至天顺}；亦见\eventref{MM-1506-EIGHT-TIGERS}{正德近侍得势}。
  \item[北元、蒙古诸部与边镇] \hyperref[register:1368-1398_hongwu]{洪武北方经营}、\hyperref[register:1399-1424_jianwen_yongle]{永乐远征}；\eventref{MM-1571-LONGQING-PEACE}{封俺答与互市}。
  \item[瓦剌、北京防务与英宗朝廷] \hyperref[register:1425-1464_hongxi_tianshun]{洪熙至天顺锚点簇}；\eventref{EM-1449-TUMU}{土木危机}。
  \item[建州、辽东与后金] \hyperref[register:1465-1505_chenghua_hongzhi]{成化、弘治边地记录}、\hyperref[register:1621-1644_tianqi_chongzhen]{天启、崇祯辽东记录}；\eventref{ML-1618-901}{1618 年辽东攻势}。
  \item[朝鲜、日本、海上与朝贡相关政权] \hyperref[register:1506-1567_zhengde_jiajing]{正德、嘉靖海防记录}、\hyperref[register:1568-1620_longqing_wanli]{隆庆、万历跨海与援朝记录}。
  \item[地方社会、流民、土司与流动作战集团] \hyperref[register:1368-1398_hongwu]{洪武地方接收}、\hyperref[register:1465-1505_chenghua_hongzhi]{成化、弘治山地行政}、\hyperref[register:1621-1644_tianqi_chongzhen]{晚明社会危机}；\eventref{MM-1465-JINGXIANG-UPRISING}{荆襄流民}与\eventref{ML-1606-901}{播州改流}。
\end{description}

\section*{主题索引}
\label{app:theme-index}

\begin{description}
  \item[建国、继承与宫廷程序] \hyperref[register:1368-1398_hongwu]{洪武}、\hyperref[register:1399-1424_jianwen_yongle]{建文--永乐}、\hyperref[register:1506-1567_zhengde_jiajing]{正德、嘉靖}；\eventref{MM-1522-GREAT-RITES-BEGIN}{大礼议}。
  \item[财政、赋役、漕运与军饷] \hyperref[register:1568-1620_longqing_wanli]{隆庆、万历}、\hyperref[register:1621-1644_tianqi_chongzhen]{天启、崇祯}；\eventref{ML-1637-901}{练饷}。
  \item[边疆、战争、卫所与道路] \hyperref[register:1399-1424_jianwen_yongle]{永乐北方经营}、\hyperref[register:1425-1464_hongxi_tianshun]{土木后重组}、\hyperref[register:1621-1644_tianqi_chongzhen]{辽东--北京危机}。
  \item[地方治理、迁徙与山地社会] \hyperref[register:1465-1505_chenghua_hongzhi]{成化、弘治}；\eventref{MM-1465-DATENG-GORGE-CAMPAIGN}{大藤峡}与\eventref{MM-1467-JIANZHOU-CAMPAIGN}{建州之役}。
  \item[港口、海防、贸易与跨海关系] \hyperref[register:1506-1567_zhengde_jiajing]{正德、嘉靖}、\hyperref[register:1568-1620_longqing_wanli]{隆庆、万历}；\eventref{MM-1567-YUEGANG-OPENING}{月港}。
  \item[教育、礼制、知识与文化] \hyperref[register:1368-1398_hongwu]{洪武官学与文书}、\hyperref[register:1506-1567_zhengde_jiajing]{嘉靖礼制}；这些记录只定位制度和文本，不把颁布倒写成普遍接受。
\end{description}

\section*{来源索引}
\label{app:source-index}

\begin{description}
  \item[中枢编年与官修汇编] \sourceref{ming-shilu}{《明实录》}；\sourceref{ming-shi}{《明史》}；\sourceref{chongzhen-changbian}{《崇祯长编》}。
  \item[制度与官员文书] \sourceref{memorials}{奏疏、题本与诏令}；\sourceref{statutes}{《明会典》及制度志}。
  \item[地方与物质材料] \sourceref{local-records}{地方志、契约、碑刻与族谱}；\sourceref{material-evidence}{遗址、器物与考古报告}。
  \item[多方材料] \sourceref{external-records}{朝鲜、越南、蒙古、日本、葡萄牙等外部记录}。
\end{description}

\section*{图件索引}
\label{app:figure-index}

\begin{description}
  \item[总体空间关系] \hyperref[fig:vol05:ming-spatial-overview]{明代国家能力的空间关系示意}。
  \item[总体时间线] \hyperref[fig:vol05:ming-chronology-timeline]{明代编年主线时间轴}。
  \item[北方与北京] \hyperref[fig:vol05:northern-frontier-beijing-grasslands]{洪武至永乐北方前沿和北京节点}。
  \item[土木与朝廷重组] \hyperref[fig:vol05:tumu-beijing-court-restoration]{土木危机、北京防务与朝廷重组}。
  \item[成化至正德的地方接口] \hyperref[fig:vol05:regional-governance-frontiers]{地方治理、边地与中介网络}。
  \item[东南海岸] \hyperref[fig:vol05:southeast-sea-coastal-networks]{防务、港口与跨海网络}。
  \item[万历财政运输] \hyperref[fig:vol05:transport-fiscal-routes]{漕运、河工与财政征解路线}。
  \item[晚明多重危机] \hyperref[fig:vol05:late-ming-crisis-multiple-fronts]{天启、崇祯的多重危机}。
  \item[辽东--北京走廊] \hyperref[fig:vol05:liaodong-beijing-corridor]{晚明辽东--山海关--北京危机走廊}。
\end{description}
